\begin{styletable}

Равенство по определению
    & $A_{\mathrm{new}} \is B$, $B \isr A_{\mathrm{new}}$
    & \verb|\is|, \verb|\isr|
\\

Равенство и его отрицание
    & $a = b, \; a \ne b$
    & \verb|=, \ne|
\\

Приближённое равенство
    & $a \approx b$
    & \verb|\approx|
\\

Умножение
    & $a \mul b$, $a \times b$
    & \verb|\mul|, \verb|\times|
\\

Неравенства строгие
    & $a < b, \; b > a$
    & \verb|<|, \verb|>|
\\

Неравенства нестрогие
    & $a \le b, \; b \ge a$
    & \verb|\le|, \verb|\ge|
\\

Инкремент
    & $1 = 0\inc$
    & \verb|1 = 0\inc|
\\

Делимость
    & $12 \divby 2$, $13 \ndivby 2$
    & \verb|\divby|, \verb|\ndivby|
\\

Плюс-минус
    & $a \pm b$, $a \mp b$
    & \verb|\pm|, \verb|\mp|
\\

Степень и индекс
    & $x^n$, $x_{i+1}$, $C_n^k$
    & \verb|x^{i+1}|, \verb|x_{i+1}|, \verb|C_n^k|
\\

Дробь
    & $\frac{a}{b}$
    & \verb|\frac{a}{b}|
\\

Корень
    & $\sqrt{x}$, $\sqrt[n]{x}$
    & \verb|\sqrt{x}|, \verb|\sqrt[n]{x}|
\\

\hline
\end{styletable}

\stylenote{
    Для умножения используем команду \texttt{\textbackslash mul} (не \texttt{\textbackslash cdot}).
    Команда \texttt{\textbackslash times} применяется для прямого произведения множеств и векторного произведения, а также в случае переноса операции умножения на новую строку.
}